\documentclass[12pt]{article}
%\documentstyle{article}
\usepackage[latin2]{inputenc}
\usepackage{amsmath}

\textwidth 14.2cm
\textheight 22.5cm
\topmargin  -0.5cm
\oddsidemargin 0.5cm
\evensidemargin 0.5cm
\pagestyle{empty}

%\partopsep -33pt
%\parsep 0pt
%\topskip -33pt
\parskip 0pt
\parindent 0pt
\setlength{\topsep}{0pt}
\setlength{\partopsep}{0pt}
\setlength{\itemsep}{-30pt}
\setlength{\parsep}{-30pt}

\makeatletter
   \def\@listi{\itemsep0pt\parsep0pt}
\makeatother




%\renewcommand{\baselinestretch}{0.9}
\renewcommand{\topfraction}{0.8}
\renewcommand{\textfraction}{0.2}
\renewcommand{\floatpagefraction}{0.8}

%\settowidth{\leftmargin}{CC}
%\addtolength{\leftmargin}{-30pt}
%\setlength{\leftmargin}{5pt}

\newcounter{denselistcounter}
\newenvironment{denselist}[1]
{ \begin{list}
  {#1{denselistcounter})}{\usecounter{denselistcounter}
  \setlength{\topsep}{-0pt}
  \setlength{\partopsep}{-0pt}
  \setlength{\itemsep}{-0pt}
  \setlength{\parsep}{-0pt}
  \setlength{\labelwidth}{6pt}
  \setlength{\labelsep}{4pt}
  \setlength{\leftmargin}{20pt}
  \setlength{\rightmargin}{20pt}
  }
}
{\end{list}}

\setlength{\unitlength}{0.1cm}
\begin{document}


\title{Gaussian folding in 2D}
\author{Discussion file: MSU -- FAU}
\date{Started: 6. September 2024; last update: 10. September 2024}
\maketitle

The aim is to express Gaussian folding 
\begin{equation}
  \psi(\mathbf{r})
  =
  \int d^3r'\, e^{-\lambda(\mathbf{r}-\mathbf{r}')^2}
  \varphi(\mathbf{r}')
\end{equation}
in cylindrical coordinates $(x,y,z)\longrightarrow(r,\phi,z)$.

\section{Direct integration in coordinate space}

One way is to perform the transformation in coordinate space.
This reads
\begin{equation}
  \psi(r,\phi,z)
  =
  \in dz'\,d\phi'\,dr'\,r'\,e^{-\lambda(z-z')^2}
  e^{(r-r')^2+2rr'(1-\cos(\phi-\phi'))}
  \varphi(r',\phi',z')
  \;.
\end{equation}
The fields $\psi$ and $\varphi$ are usually given with good azimuthal
angular momentum $m$, i.e.
\begin{equation}
  \varphi(r',\phi',z')
  =
  \varphi_m(r',z')e^{-\mathrm{i}m\phi'}
\end{equation}
and similarly for $\psi$. This reduces the transformation to
\begin{subequations}
\begin{eqnarray}
  \psi_m(r,z)
  &=&
  \in dz'\,d\eta\,dr'\,r'\,e^{-\lambda(z-z')^2} e^{(r-r')^2-2rr'}
  \varphi_m(r',z')\int d\eta e^{-2rr'\cos(\eta)-\mathrm{i}m\eta}
  \;,
\\
  \eta&=&\phi-\phi'
  \;.
\end{eqnarray}
\end{subequations}
We recall the definition of the Bessel functions
\begin{equation}
  J_n(x)
  =
  \frac{\mathrm{i}^{-n}}{2\pi}
  \int_{-\pi}^\pi d\theta\,e^{\mathrm{i}x\cos(\theta)+\mathrm{i}m\theta}
\end{equation}
which means that we are actually looking for the modified Bessel 
function
\begin{eqnarray}
  \int d\eta e^{-2rr'\cos(\eta)-\mathrm{i}m\eta}
  &=&
  \int_{-\pi}^\pi d\eta\,e^{\mathrm{i}(\mathrm{i}rr')\cos(\eta)-\mathrm{i}m\eta}
\nonumber\\
  &=&
  2\pi\mathrm{i}^n J_{|m|}(\mathrm{i}2rr')
  =
  2\pi(-)^n I_m(rr')
\end{eqnarray}
where $I_n$ are the modified Bessel functions. After all, we come to
\footnote{All steps to here need yet checking.}
\begin{equation}
  \psi_m(r,z)
  =
  (-)^m 2\pi\in dz'\,dr'\,r'\,e^{-\lambda(z-z')^2} e^{(r-r')^2-2rr'}I_m(2rr')  \varphi_m(r',z')
  \;.
\end{equation}
The direct application of this closed expression on a finite grid will
violate boundary conditions. One may ignore that if the fields are
well localized within the grid and if the folding width is not too
large. 

\section{Expansion in Laplacean eigenstates on the grid}

A better adjustment to the grid is achieved by the operator
representation as we used it in the 1D code. Starting point is
the 3D expression
\begin{equation}
  \psi(\mathbf{r})
  =
  \exp\left(\frac{\hat{\Delta}}{4\lambda}\right)
  \varphi(\mathbf{r})
\end{equation}
where $\hat{\Delta}$ is the Laplacean operator.
The Laplacean reads in cylindrical coordinates when acting
onto an eigenfunction of azimuthal angular momentum (see above)
\begin{equation}
  \hat{\Delta}
  =
  \hat{\Delta}_r-\frac{m^2}{r^2}+\partial_z^2
  \;,\;
  \hat{\Delta}_r
  =
  \partial_r^2
  +
  \frac{1}{r}\partial_r
  \;.
\end{equation}
Thus we are searching for a representation of
\begin{equation}
  \psi(\mathbf{r})
  =
  \exp\left(\frac{\hat{\Delta}_r-\frac{m^2}{r^2}+\partial_z^2}{4\lambda}\right)
  \varphi(\mathbf{r})
\label{eq:folding_axial}
\end{equation}
This raises a problem: $\hat{\Delta}_r$ is not immediately Hermitian.
This may look surprising at first glance. It becomes plausible at a
second thought. One has to take into account the volume element. To
that end, look at a typical axial matrix element
\begin{subequations}
\begin{equation}
  \langle\psi_m|\hat{\Delta}_r|\varphi_m\rangle
  =
  \int dr\,r\,\psi_m\left(\partial_r^2+\frac{1}{r}\partial_r\right)
  \varphi_m
  \;.
\end{equation}
One can show by partial integration that this is equivalent to
\begin{equation}
  \langle\psi_m|\hat{\Delta}_r|\varphi_m\rangle
  =
  \int dr\,r\,
  \left[\left(\partial_r^2+\frac{1}{r}\partial_r\right)\psi_m\right]
  \varphi_m
  \;.
\end{equation}
\end{subequations}
One can get formally rid of the volume element by considering the
scaled wavefunctions
\begin{equation}
  \tilde{\psi}_m=\sqrt{r}\,\psi_m
  \;,\;
  \tilde{\varphi}_m=\sqrt{r}\,\varphi_m
  \;.
\end{equation}
This changes the matrix element to
\begin{eqnarray}
  \langle\psi_m|\hat{\Delta}_r|\varphi_m\rangle
  &=&
  \int dr\,\tilde{\psi}_m\sqrt{r}\left(\partial_r^2+\frac{1}{r}\partial_r\right)
  \frac{1}{\sqrt{r}}\tilde{\varphi}_m
\nonumber\\
  &=&
  \int dr\,\tilde{\psi}_m\left(\partial_r^2+\frac{1}{4r^2}\right)
  \tilde{\varphi}_m
  \;.
\end{eqnarray}
The kernel $\partial_r^2+1/(4r^2)$ 
\begin{equation}
  \tilde{\Delta}_r
  =
  \sqrt{r}\left(\partial_r^2+\frac{1}{r}\partial_r\right)
  \frac{1}{\sqrt{r}}
  =
  \partial_r^2+\frac{1}{4r^2}
\end{equation}
is obviously Hermitian.

The aim is evaluate a matrix representation of the axial part of the
Gaussian folding operator in eq. (\ref{eq:folding_axial}). We call it
for compact notation
\begin{equation}
  \hat{\mathcal{G}}_r
  =
  \exp\left(\frac{\hat{\Delta}_r-\frac{m^2}{r^2}}{4\lambda}\right)
  \;.
\end{equation}
We introduce the scaled folding operator
\begin{equation}
  \tilde{\mathcal{G}}_r
  =
  \sqrt{r}\hat{\mathcal{G}}_r\frac{1}{\sqrt{r}}
  \;.
\end{equation}
We remark that
\begin{equation*}
  \tilde{\Delta}_r^n
  =
  \left(\sqrt{r}\hat{\Delta}_r\frac{1}{\sqrt{r}}\right)^n
  =
  \sqrt{r}\hat{\Delta}_r^n\frac{1}{\sqrt{r}}
  \;.
\end{equation*}
Using this auxiliary result, we can show via the Taylor expansion that
\begin{equation}
  \tilde{\mathcal{G}}_r
  =
  \exp\left(\frac{\tilde{\Delta}_r-\frac{m^2}{r^2}}{4\lambda}\right)
  \;.
\end{equation}
We first construct the matrix representation of
$\tilde{\mathcal{G}}_r$ through the eigenstates of $\tilde{\Delta}_r$
on the grid. This proceeds as
\begin{eqnarray*}
  -k_i^2\tilde{J}_m(k_ir)
  &=&
  \left(\tilde{\Delta}_r-\frac{m^2}{r^2}\right)\tilde{J}_m(k_ir)
  \;,
\\
  \tilde{\mathcal{G}}_r(r,r')
  &=&
  \sum_i\tilde{J}_m(k_ir)\exp\left(-\frac{k_i^2}{4\lambda}\right)
  \tilde{J}_m(k_ir')
\end{eqnarray*}
where $\tilde{J}_m(k_ir)=\sqrt{r}{J}_m(k_ir)$ represents the Bessel
functions on the grid.
The folding operator kernel is finally 
\begin{equation}
  {\mathcal{G}}_r(r,r')
  =
  \frac{1}{\sqrt{r}}\tilde{\mathcal{G}}_r(r,r')\sqrt{r}
  \;.
\end{equation}
\centerline{\em ... boundary $r\rightarrow 0$, $r\rightarrow
  R_\mathrm{box}$ yet to be discussed ...}

\section{Expansion in Bessel functions}

We come back to the idea of using the analytically given eigenstates 
of the axial Laplacean, namely the Bessel function $J_m(x)$,
directly. The Bessel functions are given as solution of
\begin{equation}
  \left(\hat{\Delta}_r-\frac{m^2}{r^2}\right)J_m(kr)
  =
  \left(
  \partial_r^2
  +
  \frac{1}{r}\partial_r
  -
  \frac{m^2}{r^2}
  \right)J_m(kr)
  =
  -k^2J_m(kr)
\end{equation}
where $J_{-m}=J_m$, i.e. the sign of $m$ plays no role.
The Bessel functions to different $k_n$ are orthogonal to each other.
But they are not normalized. To use them as expansion basis, one has 
first to normalize the $J_m$.

Not all $k$ are possible in a finite numerical box with radius
$R_\mathrm{box}$. Two kind of boundary conditions are used
\begin{subequations}
\begin{eqnarray}
  J_m(k_nR_\mathrm{box})
  &=&
  0
  \;,
\\
  J'_m(k'_nR_\mathrm{box})
  &=&
  0
  \;.
\end{eqnarray}
\end{subequations}
Both yield a series of discrete values $k_n$, or $k'_n$ respectively.
The series is terminated from above due to finite resolution on the
grid by spatial discretization. 
After all, we obtain two expansions sets $\left\{J_m(k_nr)\right\}$ and
$\left\{J_m(k'_nr)\right\}$ depending on the wanted boundary conditions.

\section{Testing through simple expectation values}

For a test, we take up the analytical formula from the case of
spherical 1D (see the discussion file \texttt{doc/test\_folding.pdf}
in the 1D repository). Given a s.p. function in spherical representation
\begin{subequations}
\begin{equation}
  \varphi(r_\mathrm{sph},\vartheta,\phi)
  =
  r_\mathrm{sph}^{l}Y_{ll}(\Omega)
  \exp\left(-\frac{r_\mathrm{sph}^2}{\delta r_\mathrm{sph}^2}\right)
\end{equation}
yields the folded function
\begin{equation}
  \tilde{\varphi}(r_\mathrm{sph},\vartheta,\phi)
  =
  (-)^l\left(\frac{\delta r_\mathrm{sph}^2}{\Delta r^2+\delta r^2}\right)^{l+3/2}
  r_\mathrm{sph}^{l}Y_{ll}(\Omega)
  \exp\left(-\frac{r_\mathrm{sph}^2}{\Delta r^2+\delta r^2}\right)
  \;.
\end{equation}
\end{subequations}
This relation will now be reshaped into axial coordinates
$(r,\Phi,z)$. The spherical harmonics $Y_{ll}$ is independent of $z$
Thus the relation becomes
\begin{subequations}
\begin{eqnarray}
  \varphi(r,\phi,z)
  &=&
  r_\mathrm{sph}^{l}e^{\mathrm{i}m\phi}
  \exp\left(-\frac{r_\mathrm{sph}^2+z^2}{\delta r^2}\right)
  \;,
\\
  \tilde{\varphi}(r,\phi,z)
  &=&
  (-)^l\left(\frac{\delta r_\mathrm{sph}^2}{\Delta r^2+\delta r^2}\right)^{l+3/2}
  r_\mathrm{sph}^{l}e^{\mathrm{i}m\phi}
  \exp\left(-\frac{r_\mathrm{sph}^2+z^2}{\Delta r^2+\delta r^2}\right)
  \;.
\end{eqnarray}
\end{subequations}
To test folding in $z$ direction, we consider the $z^2$ expectation
value which is
\begin{subequations}
\begin{equation}
  \langle z^2\rangle_\varphi
  =
  \frac{\int dz d\phi dr\,r\,e^{-\mathrm{i}m\phi}\varphi(r,\phi,z)z^2}
       {\int dz d\phi dr\,r\,e^{-\mathrm{i}m\phi}\varphi(r,\phi,z)}
  =
  \frac{\delta r^2}{2}
\end{equation}
and for the folded test function
\begin{equation}
  \langle z^2\rangle{\tilde\varphi}
  =
  \frac{\delta r^2+\Delta r^2}{2}
  \;.
\end{equation}
\end{subequations}
The change through $z$-folding is thus  $\Delta r^2/2$.

For the $r$ direction, we consider the expectation value
\begin{equation}
  \langle r^2\rangle_\varphi
  =
  \frac{\int dz d\phi dr\,r\,e^{-\mathrm{i}m\phi}\varphi(r,\phi,z)r^2}
       {\int dz d\phi dr\,r\,e^{-\mathrm{i}m\phi}\varphi(r,\phi,z)}
  \;.
\end{equation}
Using the integrals
\begin{eqnarray*}
  \int_0^\infty dx\,x^{2k}\,e^{-\lambda x^2}
  &=&
  \frac{1}{2}\,\sqrt{\frac{\pi}{\lambda}}\,\frac{(2k-1)!!}{2^k\lambda^k}
\\
  \int_0^\infty dx\,x^{2k+1}\,e^{-\lambda x^2}
  &=&
  \frac{1}{2}\,\frac{k!}{\lambda^{k+1}}
\end{eqnarray*}
we derive
\begin{equation}
  \langle r^2\rangle_\varphi
  =
  \frac{l+2}{2}\delta r^2
  \;,\;
  \langle r^2\rangle_{\tilde\varphi}
  =
  \frac{l+2}{2}(\delta r^2+\Delta r^2)
  \;.
\end{equation}
These are the tests used in subroutine \texttt{test\_Gaussfolding} in
the file \texttt{m\_grid.f90}.


\end{document}

\begin{equation}
\end{equation}
\begin{equation}
\end{equation}
\begin{equation}
\end{equation}
